(1) 一段实验室时间 \(dt\) 内光帆只迎来相对长度 \((c-v)dt\) 内的光，截获能量为 \(P_0(1-\beta)dt\)。反射光子的实验室能量比 \(r=(1-\beta)/(1+\beta)\)，反向传播使动量转移为 \(d p_{\rm sail}=(dE_i+dE_r)/c\)。因此
\[
F(v)=\frac{P_0(1-\beta)}c(1+r)
=\frac{2P_0}{c}\frac{1-\beta}{1+\beta}.
\]
\(v=0\) 时 \(F=2P_0/c\)。
(2) \(\gamma\beta=\sinh\chi,\ \beta=\tanh\chi,\ (1-\beta)/(1+\beta)=e^{-2\chi}\)。所以
\[
mc\cosh\chi\,\dot\chi=(2P_0/c)e^{-2\chi},
\qquad \frac{dt}{d\chi}=\frac{mc^2}{2P_0}e^{2\chi}\cosh\chi.
\]
取 \(t(0)=0\)，积分
\[
t(\chi)=\frac{mc^2}{4P_0}\left[\frac{e^{3\chi}-1}{3}+e^\chi-1\right].
\]
同理 \(dx/d\chi=c\tanh\chi\,dt/d\chi=mc^3e^{2\chi}\sinh\chi/(2P_0)\)，得
\[
x(\chi)=\frac{mc^3}{4P_0}\left[\frac{e^{3\chi}-1}{3}-(e^\chi-1)\right].
\]
两式之差 \(ct-x=\frac{mc^3}{2P_0}(e^\chi-1)\ge0\)。起动加速度 \(\dot v(0)=2P_0/(mc)\)。
(3) 大快度下 \(t\sim mc^2e^{3\chi}/(12P_0)\)，从而
\[
e^\chi\sim\left(\frac{12P_0t}{mc^2}\right)^{1/3},\
\gamma\sim\frac12\left(\frac{12P_0t}{mc^2}\right)^{1/3},\
1-\beta\sim2\left(\frac{12P_0t}{mc^2}\right)^{-2/3}.
\]
任何有限 \(t\) 都有 \(x<ct\)。长期完美反射、无限平行光束、无热损和恒定 \(P_0\) 是理想化前提；真实光束衍射及光帆面积会限制加速。
(4) 截获入射功率 \(P_i=P_0(1-\beta)\)，反射功率 \(P_r=P_i(1-\beta)/(1+\beta)\)。其差
\[
P_i-P_r=P_0(1-\beta)\frac{2\beta}{1+\beta}
=Fv=\frac d{dt}(\gamma mc^2),
\]
恰为光帆动能增加功率。以截获功率为分母，转化效率为 \(2\beta/(1+\beta)\)，在 \(\beta\to1^-\) 时趋一；但实际截获功率同时趋零，因而没有无限功率或超光速问题。连续流模型对应很多低反冲光子，不可拿单光子动能守恒替代整体的动量通量账本。
